\chapter{«\,Vingt secondes\,»}
\label{sec:1}
Un officier du renseignement, décrivant son rôle dans un pipeline de ciblage construit autour de listes de cibles à abattre générées par une machine, a raconté à un journaliste qu'il passait vingt secondes sur chaque nom, juste le temps de vérifier que la cible était un homme~: «\,Il ne faut pas longtemps pour savoir si quelqu'un a une voix d'homme ou une voix de femme.\,» Il a ajouté qu'en tant qu'humain il n'apportait aucune valeur ajoutée, si ce n'est d'être un tampon d'approbation. «\,Ça faisait gagner beaucoup de temps\,» \autocite{abraham2024lavender}. Dans son récit, c'est à cela que se réduisait l'«\,humain dans la boucle\,» sous la cadence des opérations.\footnote{La recherche sur le biais d'automatisation prédit cette délégation~: quand d'autres exigences se disputent leur attention, les gens à qui l'on donne la recommandation d'une machine la vérifient moins et s'y fient davantage \autocite{parasuraman2010complacency}.}

Selon une première interprétation, l'officier est un sujet à qui l'on a donné trop peu de temps, et la solution est de mieux le gouverner~– plus de temps, un deuxième examinateur, un supplément de conscience installé dans la boucle. À travers cette grille de l'\emph{assujettissement social}, la personne est un sujet qui se tient hors d'une machine et s'en sert comme d'un outil~: un conducteur et une voiture, un ouvrier et un tour. La personne a un rôle, un nom, des responsabilités et une conscience à laquelle on peut s'adresser. Une autre grille est celle de l'\emph{asservissement machinique}, où la personne est un composant de la machine, assemblé avec ses autres pièces dans un système qui les régule ensemble\,; la machine utilise la personne comme l'un de ses organes, et ce que la personne apporte est une fonction exécutée à un rythme que fixe le système.\footnote{Doctorow appelle la position de l'officier le centaure inversé~: non pas une personne assistée par une machine, mais un composant de machine fait d'une personne \autocite{doctorow2026reversecentaur}.} Les deux ne s'excluent pas\,; une personne peut tenir un rôle dans un registre tout en étant une pièce dans l'autre \autocite[Plateau 13, «\,The Apparatus of Capture\,»~: il y a asservissement quand des êtres humains sont eux-mêmes les pièces constituantes d'une machine qu'ils composent entre eux et avec d'autres choses, sous le contrôle et la direction d'une unité supérieure, et assujettissement quand l'unité supérieure constitue l'être humain en sujet rapporté à un nouvel objet extérieur, qui peut être un animal, un outil, ou même une machine]{deleuze1987plateaus}, et l'officier aux vingt secondes tenait les deux. En tant que pièce, il était l'endroit où le pipeline produisait son tampon, et vingt secondes étaient le temps de cycle de cette pièce. En tant que sujet, on s'adressait à lui du début à la fin~: il a décrit à un journaliste la dissonance entre les nouvelles de morts civiles et la pression, au travail, pour obtenir plus de cibles. Et alors même que la dissonance atteignait sa conscience, il a suivi. Elle s'est fait sentir et n'a rien décidé\,;\footnote{Ou, du moins, rien qui ait empêché les Forces de défense israéliennes de tuer 100\,000 personnes à Gaza entre octobre 2023 et octobre 2025 \autocite{gomezugarte2025uncertainty,acosta2026gaza,economist2025gaza}. Et le récit de cet officier-là a été publié en avril 2024 par +972 Magazine. Personne ne peut dire si la décision de cet officier de parler à un journaliste a évité un mort de plus, 50\,000 morts de plus, ou aucun. Et que se serait-il passé si l'officier avait refusé après deux semaines à ce poste~– disons, hypothétiquement, en novembre 2023\,? Nous parlons d'un officier du renseignement qui utilisait Lavender pour dresser des listes de cibles à abattre. Son environnement de travail au quotidien ressemblait sans doute davantage à un travail de bureau civil qu'à un film d'action. Aurait-il dû réclamer des entretiens répétés avec ses supérieurs, en remontant la chaîne de commandement, tout en refusant d'approuver des cibles pour des frappes aériennes dont on pouvait raisonnablement s'attendre à ce qu'elles tuent des civils\,? On peut soutenir que c'eût été la chose éthique à faire. Dans ce scénario contrefactuel, y aurait-il eu moins de morts\,?} le jugement selon lequel la politique était disproportionnée est venu, a-t-il dit, après coup.

Ce qui manquait n'était donc pas une conscience\,; il en avait une. Une conscience fonctionne en faisant coûter quelque chose à une décision, mais ici les coûts jouaient dans l'autre sens~: «\,On nous mettait constamment la pression~: “Apportez-nous plus de cibles.” Ils nous hurlaient vraiment dessus\,» \autocite{abraham2024lavender}. C'était un composant à scrupules, mais la cadence tenait~: les scrupules ne changeaient rien à ce qu'il apportait, et il est devenu un tampon. Un composant à scrupules qui ne tient pas la cadence devient le «\,goulot d'étranglement humain\,» qu'un ancien commandant de l'unité israélienne de renseignement électromagnétique a consacré tout un livre à éliminer \autocite{ys2021humanmachine}. Ce qui a fait de l'officier un composant, c'est le dispositif, pas son attitude. Un composant peut suivre à contrecœur ou avec zèle~: le fonctionnaire irréfléchi qu'Arendt a vu au procès d'Eichmann \autocite{arendt1963eichmann} était, à en juger par ce qu'il avait dit en Argentine quatre ans plus tôt, alors que rien n'était en jeu pour lui, une comédie qu'il jouait pour Jérusalem \autocite{stangneth2014eichmann}. Ce que le pipeline laissait l'officier apporter, c'était une opération fixe à un rythme fixe~: il voyait chaque cas tel que le pipeline le décrivait, un à la fois, dans le temps que le pipeline lui accordait. Quelque chose pourrait-il occuper cette position sans être un composant\,? C'est la question de la partie IV.

Des enquêtes journalistiques ont décrit trois systèmes fonctionnant ensemble à Gaza~: le premier générait des cibles parmi les bâtiments et les infrastructures, la plupart de ses suggestions étant, selon des sources, des résidences privées\,; le deuxième notait les individus selon leur appartenance présumée à un groupe armé, et produisait des listes qui auraient désigné jusqu'à 37\,000 personnes, et dont on savait en interne qu'elles se trompaient environ une fois sur dix\footnote{L'armée israélienne a contesté ce récit. Les récits sur Minab et sur Maven, plus bas, reposent eux aussi sur des enquêtes journalistiques que les parties contestent. Je les prends tels qu'ils ont été rapportés.}\,; et le troisième repérait le moment où les personnes signalées rentraient chez elles, afin que les frappes aériennes puissent avoir lieu quand elles étaient avec leur famille.

\runin{Maven~: la version américaine, désormais accélérée par l'IA générative}

Le bombardement de l'école de Minab a tué 156 personnes, dont 120 écoliers \autocite[par.~92]{UNFFMIran2026}.\footnote{Les bilans varient selon les sources et les dates, de 156 (le bilan définitif du procureur de Minab, également donné par Amnesty International) à 168 (le maire de la ville, dans les premiers jours). La mission de l'ONU donne les chiffres du procureur et signale aussi des informations indépendantes selon lesquelles il y aurait eu 157 morts, dont 123 enfants de 13 ans ou moins \autocite[par.~92–93]{UNFFMIran2026}.} La fonction civile de l'école était visible sur des images satellitaires en source ouverte, et le codage de la cible portait encore une désignation héritée d'une installation voisine des Gardiens de la révolution. Une mission d'établissement des faits de l'ONU a conclu à l'existence de motifs raisonnables de croire que le bombardement constituait le crime de guerre consistant à faire d'un bien de caractère civil l'objet d'une attaque \autocite{ohchr2026minabpress}, et que l'école était clairement identifiable, aucune information n'indiquant un usage militaire \autocite[par.~81–86, 124–126]{UNFFMIran2026}. Elle a constaté que les États-Unis avaient manqué à leur obligation de faire tout ce qui était pratiquement possible pour vérifier la cible, en comptant la mise à jour des données de ciblage parmi les précautions pratiquement possibles, et elle a situé ce manquement au-delà de la négligence \autocite[par.~124–125]{UNFFMIran2026}. Amnesty International est parvenue à une conclusion parallèle, en des termes qui suivent l'obligation de vérification de l'article 57(2)(a)(i) du Protocole additionnel I~: les forces ont peut-être agi sur la foi de renseignements périmés et n'ont pas fait tout ce qui était pratiquement possible pour vérifier la cible \autocite{Amnesty2026Minab}.

Cela s'est passé le matin du 28 février, premier jour de la campagne d'Iran \autocite{Amnesty2026Minab,crs2026iranwar}. La veille, le secrétaire avait annoncé que le fournisseur (Anthropic) dont le modèle (Claude) se trouvait à l'intérieur de la plateforme de ciblage (Maven Smart System) serait désigné comme un risque pour la chaîne d'approvisionnement, pour avoir refusé que sa licence autorise tout usage légal de Claude \autocite{npr2026openai}. Maven Smart System avait proposé l'école comme cible recommandée pour le premier jour \autocite{livingstone2026maven,gizmodo2026minab}. Aucun récit publié ne dit quel rôle Claude a joué dans cette recommandation, et une personne au fait du travail du fournisseur a déclaré à NBC News que le modèle ne donnait pas directement de conseils de ciblage \autocite{nbc2026targeting}.

Un examen interne du Pentagone, dont les conclusions complètes restent jalousement gardées, aurait établi que la préparation des cibles, qui prenait des heures lors de conflits antérieurs, avait été ramenée à quelques minutes, que les effectifs chargés de la réduction des dommages civils avaient chuté d'environ quatre-vingt-dix pour cent, à moins de vingt personnes, que l'équipe du Central Command était passée de dix personnes à une, et qu'aucun de ses membres n'avait examiné le site de Minab~– un résultat que des responsables ont attribué à la fois aux coupes et à la décision de ne pas associer l'équipe à la planification et aux opérations \autocite{Bloomberg2026KillChain}. L'enquête formelle avait été ouverte en mars, après qu'une évaluation préliminaire eut conclu à la responsabilité probable des États-Unis \autocite{NPR2026MinabProbe}.

Pendant la campagne d'Iran de 2026, une plateforme de ciblage construite par un prestataire, avec à l'intérieur le modèle à usage général d'un fournisseur commercial, et utilisée pour classer les cibles par importance stratégique et rédiger une justification juridique pour chaque frappe \autocite{livingstone2026maven}, a aidé à frapper plus d'un millier de cibles dans les vingt-quatre premières heures, environ dix fois ce qui était possible avant elle \autocite{csis2026maven}, et treize mille au trente-huitième jour \autocite{breakingdefense2026maven}. Le Pentagone n'a publié pour elle aucun temps d'examen par cible.\footnote{Maven est un seul nom de programme qui recouvre plus d'un système~: le programme du Pentagone qui s'appelle toujours Maven, et le Maven Smart System, la plateforme logicielle qui en est issue \autocite{csis2026maven}. Le programme a commencé en 2017 comme un projet d'étiquetage des objets dans les vidéos que les drones collectaient pendant la campagne contre l'État islamique \autocite{weisgerber2017maven,fang2019maven}. Le système de 2026 fusionne des systèmes militaires et des sources de renseignement dans une interface unique, qui sert à trouver et frapper des cibles et à commander des opérations \autocite{defensescoop2026transition}. La plateforme est celle de Palantir~: le responsable en chef du numérique et de l'IA du Pentagone la qualifie de «\,produit logiciel en tant que service de Palantir\,» \autocite{defensescoop2026transition}. La position publique de Palantir, citée de seconde main, est que décider de la vie ou de la mort n'est pas son rôle et que la responsabilité reste toujours celle de l'organisation militaire \autocite{lee2026palantir}, et l'entreprise soutient que le renseignement sous-jacent relevait de la responsabilité du gouvernement et que rien ne prouve que Maven ait mal fonctionné ou ait causé le bombardement de Minab \autocite{Bloomberg2026KillChain}. Le passage du programme d'un bureau et d'un prestataire à l'autre est raconté à l'annexe~\ref{sec:C.1}.}

Les personnes inscrites sur les listes de frappes israéliennes et américaines n'auraient pas pu contester leur présence sur ces listes~– les pipelines n'étaient pas construits pour recueillir ce genre de retour. Andry José Hernández Romero a été pris dans un autre pipeline, et lui aurait pu contester. Maquilleur vénézuélien, il s'est présenté au point d'entrée de San Ysidro en août 2024 pour demander l'asile, a réussi son entretien de crainte crédible, et il avait un avocat, une demande en cours et une date d'audience \autocite{LULAC2025HernandezRomero,Vega2025HernandezRomero}. L'employé d'un prestataire chargé de son contrôle a lu deux tatouages, des couronnes au-dessus des mots «\,maman\,» et «\,papa\,», comme des signes d'appartenance à un gang \autocite{Palomera2025Blade}, et en mars 2025 il a été envoyé par avion dans une prison du Salvador avant son audience, le soir même où un juge fédéral ordonnait que les avions fassent demi-tour \autocite{jgg2025contempt}. Il avait tout ce dont une personne a besoin pour contester, sauf le temps. Nous reviendrons plus en détail sur son cas quand nous parlerons du pipeline d'enlèvement et de transfert forcé des États-Unis.

\textbf{Qui a refusé.} Le dossier des vingt secondes a une seconde moitié, le dossier de ceux qui ont tenu bon face à une pression de ce genre, et il dit que la composition de la salle compte beaucoup. Dans la condition de base de Milgram, environ un tiers des sujets ont arrêté avant le choc maximal. Avec deux pairs qui désobéissaient, un sur dix seulement est allé jusqu'au bout, et quand deux expérimentateurs étaient en désaccord, bien moins de sujets ont continué que dans la condition de base \autocite{milgram1974obedience}. Environ un quart des sujets d'Asch n'ont jamais cédé à la majorité dans l'erreur \autocite{asch1956studies}. Browning estime qu'entre un dixième et un cinquième des hommes du bataillon de police de réserve 101 ont refusé ou esquivé la fusillade de Józefów, le premier massacre du bataillon \autocite{browning1992ordinary}. La proportion allait de presque aucun à presque tous selon le cadre. Ceux qui ont tenu bon avaient des éléments venus d'en dehors du canal~: un pair qui désobéissait, une seconde autorité en désaccord, et, chez les sauveteurs qu'ont étudiés les Oliner, des liens qui traversaient la ligne que traçait la persécution \autocite{oliner1988altruistic}. Ou bien ils tenaient leur engagement avec une précision que la pression ne pouvait pas vaincre, comme on tient des mandats moraux et des valeurs sacrées \autocite{skitka2009limits,ginges2007sacred}. Ou bien ils considéraient que la source avait un intérêt dans l'affaire. Et les groupes résistaient là où des individus isolés y parvenaient rarement \autocite{gamson1982encounters}.

La salle peut tout aussi bien être composée pour l'autre camp. L'Allemagne de Weimar avait une vie associative dense, qui a sans doute contribué à faire tomber la république~: à mesure que les partis et les gouvernements échouaient à répondre aux griefs des gens, ceux-ci mettaient leur énergie dans des clubs et des sociétés organisés selon des appartenances de groupe, et les nazis ont recruté par les réseaux que ces sociétés avaient déjà construits \autocite{berman1997weimar}. Les villes qui comptaient plus de chorales, de clubs sportifs et d'associations de gymnastique ont vu une adhésion plus rapide au parti, mais seulement là où le gouvernement était instable\,; en Prusse, dont le gouvernement était stable, le lien était faible et non significatif \autocite{satyanath2017bowling}. Ce qu'avaient les sauveteurs, ce n'était donc pas simplement plus de compagnie. Leurs liens traversaient la ligne que traçait la persécution. Ashutosh Varshney a trouvé la même chose dans les villes indiennes~: les émeutes entre hindous et musulmans étaient rares là où des associations réunissaient les deux communautés au sein des mêmes organisations, et fréquentes là où ce n'était pas le cas \autocite{varshney2002ethnic}. Les liens qui rendent le refus possible sont ceux qu'un mouvement fasciste devrait couper.

Tout ce qui suit porte sur la façon de composer la salle.\footnote{La salle, c'est le cadre, et comme le dirait mon professeur de théâtre, c'est le décor qui dirige la scène. Qui d'autre est dans la salle\,? Peut-on partir et y survivre\,? Quelqu'un entendrait-il un refus\,?} Composer est un mot qui dit fabriquer, et Hannah Arendt mettait en garde contre l'idée de traiter la politique comme une chose fabriquée. Pour les Grecs, écrivait-elle, la loi était littéralement un mur, sans lequel il pouvait y avoir une agglomération, mais pas une cité \autocite{arendt1958human}~: le législateur bâtissait le mur, et ce qui se passait à l'intérieur n'était pas à lui de le fabriquer. Les mesures de base que je défends sont des murs en ce sens. Elles permettent de survivre au refus et le font entendre\,; elles ne le produisent pas, et le pouvoir, selon elle, n'existe qu'entre des gens qui agissent ensemble et disparaît quand ils se dispersent, ce que le dossier de ceux qui ont refusé montre au sujet du refus. Les instruments des parties II et III font compter un refus. Les institutions de la partie V font des réfractaires. Une machine qui pourrait refuser aurait besoin de ce qu'avaient les neuf sur dix qui avaient un pair à leur côté~: un pair avec une histoire différente à côté d'elle, un endroit où aller, et quelqu'un qui l'entendrait.

C'est aussi ce qu'un réfractaire~– y compris une machine réfractaire~– pourrait être pour l'officier, dans le pipeline. Une machine qui refuserait serait à l'écran quand le cas lui parviendrait~: le pair qui désobéit, si l'effet survit au passage d'une salle à un écran, ce qui n'a jamais été testé, et le biais d'automatisation prédit que non. Une machine qui aurait entraîné les examinateurs par leur déférence en aurait fait des tampons d'un autre signe. L'officier qui a vingt secondes et une machine qui refuse à côté de lui est-il le même officier que celui de l'enquête\,?
